\section{Nonnegative opening: determinant and trace faces}
\label{sec:nonnegative-opening}
\label{subsec:deg-nonnegative}

The high-height estimate reduces the remaining problem to bounded
rescaled heights. At their lower boundary, the determinant chart has a
hyperbolic saddle, while the trace chart includes its saddle--node limit.
The determinant count requires the absolute first variation, including
both tail gains. The trace count also requires the additional outgoing
saddle. We treat
both passages in the affine input and with the divided closing of
Lemma~\ref{lem:closing}.

\begin{lemma}[Finite coefficient--phase atlas]
\label{deg:lem-positive-atlas}
Fix the signed high-height constant \(K_H\) and put \(K=2K_H\).
For \(\alpha_0\ge0\), below \(z(\gamma)\le K(\alpha_0+\alpha_1^2)\),
the nonzero coefficients have the overlapping descriptions
\begin{align}
 &\alpha_0=r_D^2>0,\qquad \alpha_1=r_DA_D,\qquad |A_D|\le P_D,
 \label{deg:determinant-chart}\\
 &\alpha_0=r_L^2a_L,\qquad \alpha_1=\sigma r_L,\qquad
 0\le a_L\le a_*,\quad \sigma\in\{-1,1\},
 \label{deg:trace-chart}
\end{align}
where \(P_D>2/\sqrt{a_*}\). The determinant and signed-trace
receivers below cover every low-height through input, including regular
continuations, by overlapping stopped and compact strict intervals.
Their cuts precede a displacement zero. If \(\alpha_0=\alpha_1=0\),
every positive-height through input is in the signed high-height range.
\end{lemma}

\begin{proof}
Use \eqref{deg:determinant-chart} when
\(|\alpha_1|\le P_D\sqrt{\alpha_0}\); otherwise set
\(r_L=|\alpha_1|\), so \(a_L<a_*/4\). The face
\(\alpha_0=0\), \(\alpha_1\ne0\), has \(a_L=0\).
These are whole-coefficient alternatives with a strict overlap and no
restriction on either \(\gamma/r_D\) or \(\gamma/r_L\).
Their bounded normalized heights are respectively \(K(1+P_D^2)\)
and \(K(1+a_*)\). Their physical phase covers are constructed in
Lemma~\ref{deg:lem-determinant-preparation} and
Proposition~\ref{deg:prop-stopped-trace}. For the trace limiting field
\(X'=a_L+\sigma X-Y\), \(Y'=-XY\), at \(a_L=0\) its centre
manifold has \(Y=\sigma X+O(X^2)\), hence
\(\dot Y=-\sigma X^2+O(X^3)\); its weak quadratic coefficient is
nonzero. The double-zero coefficient face has no positive low-height
input, and is covered by Corollary~\ref{rew:cor-high-derivative-signed}.
\end{proof}

\subsection{Preparation on the actual determinant sections}

\begin{lemma}[Determinant section preparation]
\label{deg:lem-determinant-preparation}
Fix \(P_D<\infty\), the passive compact coefficient block, and the
original buffered windows. For independent sufficiently small
\(r=r_D>0\) and \(|\gamma|\), \(\alpha_0=r^2\),
\(\alpha_1=rp\), \(|p|\le P_D\), choose a fixed positive unscaled
raw-height cap \(z_*\). On the actual interval
\begin{equation}\label{deg:det-actual-receiver}
 \mathcal I_\theta=\{x:a_\theta+x\in J_\theta^{\rm out},\quad
       0<H_\theta(x)<z_*\},\qquad
 H_\theta(x)=z_\theta(\gamma;a_\theta+x),
\end{equation}
the stopped transition has the exact composition
\eqref{sharp:upper-det-composition}, with the two absolute positive
gains \eqref{sharp:upper-det-scales} and actual cut gap
\eqref{deg:det-actual-gap}. The normalized end maps and their inverses
have common ordinary \(C^3\) bounds and derivative units. The incoming
and outgoing local chords are independent. The compact-spectrum saddle
estimate \eqref{sharp:upper-det-dulac} and the complete absolute
variation \eqref{sharp:upper-det-absolute} hold, including saddle ratio
one. Every phase derivative fixes all coefficients and the actual \(m\).

The low stopped and compact strict receivers, their full regular
continuations, and the signed high receiver cover \(\mathcal I_\theta\)
as in \eqref{deg:det-height-cover}. On the strict and high pieces the
full \(\Pi'\), including divided closing, is uniformly bounded.
Retain intervening inputs, equality dividers and all original clips;
exclude natural stopped endpoints. The parameters \(m,r,\gamma\)
and the two local chord lengths are independent.
\end{lemma}

\begin{proof}
We construct the two section maps with their absolute gains before
estimating the stopped derivative. Keep \(x=h-a_\theta\), \(\Pi=\mathcal C_{\theta,m}\circ
T_\theta(a_\theta+\cdot)\), and all original finite-side factors.
The original domain is
\[
 J_\theta^{\rm out}=\mathcal T_\theta\cap\operatorname{int}K_+^\sharp
 \cap T_\theta^{-1}(\operatorname{int}K_-^\sharp)
 \cap\{\sigma_\theta<d_{\rm out}\}.
\]
Order of the raw passage and central height makes
\eqref{deg:det-actual-receiver} an interval, before capture is invoked.
With \(v=\gamma+rX\), \(z=r^2Y\), a positive time scaling gives
\begin{equation}\label{deg:det-field}
 \begin{split}
 X'&=1+pX-3\gamma X^2-rX^3-\ell(\gamma+rX)Y,\\
 Y'&=-Y\{q_\gamma X+rX^2+Cr(\gamma+rX)Y\},\qquad q_\gamma=1+2\gamma.
 \end{split}
\end{equation}
At \(r=\gamma=0\), the saddle \((0,1)\) has eigenvalues
\((p\pm\sqrt{p^2+4})/2\). Their uniform nonzero moduli give common
saddle boxes, transverse cuts \(X=\pm L_c\), and persistent branch
intersections for independent small \(r,\gamma\).
The incoming gap is exactly
\begin{equation}\label{deg:det-actual-gap}
 \xi_D=Y(L_c)-Y_b^+(L_c)>0.
\end{equation}
For a limiting through graph starting at \(Y(0)>1\), put
\(D=Y-1-pX\). The identity \(\dot D|_{D=0}=-XY\) excludes a
first exit from \(D>0\) on the negative half-axis; reverse time treats
the positive half-axis. Thus every such graph is strict with quadratic
tails. The gap is continuous, strictly increasing, and tends uniformly
to zero at the stopped boundary. A fixed larger strict graph bounds the
cut heights by order. On any bounded central-height range, the complement
\(\xi_D\ge\xi_*/2\) therefore has a finite compact strict cover with
positive finite-cut height and speed margins.

Choose enlarged tail cuts \(\pm L_t\) and buffered auxiliary ports
with \(0<d_\epsilon<d_+\), joined exactly to the original sections.
For \(X=\epsilon R\), \(\epsilon=\pm1\), the outward equation is
exactly \eqref{eq:parabolic-tail} with \((a,s)=(1,r)\):
\begin{equation}\label{deg:det-tail-graph}
 Y_R=F^\epsilon(R,Y):=F_\epsilon(R,Y)|_{a=1,s=r}.
\end{equation}
The enlarged corridor has \(Y\) and its denominator comparable to
\(R^2\) up to \(d_\epsilon/r<d_+/r\). Choose
\(|\lambda_\gamma|<\nu_0<1\). The intrinsic tail estimates with
reference graph \(Y_*\) and first variation \(j\) give
\begin{equation}\label{deg:det-tail-action}
 F_Y=\lambda_\gamma/R+O(r+|p|R^{-2}+R^{-3+\nu_0}),\qquad
 j=(R/L_t)^{\lambda_\gamma}U_\epsilon,\quad U_\epsilon^{\pm1}\le C,
\end{equation}
\begin{equation}\label{deg:det-rational-derivatives}
 |F_{Y^h}|\le C_h\{rR^{-2h+2}+|\gamma|R^{-2h+1}
              +|p|R^{-2h}+R^{-2h-1}\},\qquad h=2,3.
\end{equation}
The first remainder has bounded integral to the actual endpoint.
Use the normalization of Lemma~\ref{lem:scalar-normalization},
\begin{equation}\label{deg:det-normalized-tail}
 \eta=(Y-Y_*)/j,\qquad
 \eta_R=N=\{F(R,Y_*+j\eta)-F(R,Y_*)-jF_Y(R,Y_*)\eta\}/j.
\end{equation}
Its changed kernel weight is
\begin{equation}\label{deg:det-normalized-kernel}
 \|N\|_{C^3}\le C\{rR^{-2+\nu_0}+|\gamma|R^{-3+\nu_0}
                      +P_DR^{-4+\nu_0}+R^{-5+\nu_0}\}.
\end{equation}
Choose \(L_t\) large, then shrink \(r,|\gamma|\) independently.
The lemma with \(k=2\) supplies enlarged chords and smaller inverse
images, remaining in the corridor, with
\begin{equation}\label{deg:det-phase-variations}
 \|\mathcal T_\epsilon\|_{C^3}
 +\|\mathcal T_\epsilon^{-1}\|_{C^3}\le C,\qquad
 0<c\le\mathcal T_\epsilon'\le C.
\end{equation}

At the actual port the raw centre and gain are
\(z_{\rm ref}=r^2Y_*(d_\epsilon/r)\), \(d=r^2j(d_\epsilon/r)>0\), with
\(d\le Cr^{2-\nu_0}\). The fixed-port proof of
Lemma~\ref{ann:p0-fixed-tail-ports} admits the whole raw images:
choose the auxiliary abscissa ratio from the window reserve, then shrink
the coefficients on the remaining compact bridges. That proof uses only
the common quadratic corridors just established. On these images
Lemma~\ref{phys:lem-c3-adapters}
gives the exact ordinary \(C^3\) normalization
\begin{equation}\label{deg:det-adapter-normalization}
 \widehat G(u)=\frac{G(z_{\rm ref}+du)-G(z_{\rm ref})}
 {G'(z_{\rm ref})d}=u+O_{C^3}(d),
\end{equation}
with common inverse bounds. Its divided-closing rows
\eqref{phys:eq-closing-jets} give, including \(m=0\),
\begin{equation}\label{deg:det-divided-closing}
 \|\mathcal C_{\theta,m}\|_{C^3}\le C,
 \qquad 0<c\le\mathcal C_{\theta,m}'\le C.
\end{equation}

Finite composition, with the positive-tail inverse on the incoming
side, now gives the two exact absolute embedding gains
\begin{equation}\label{sharp:upper-det-scales}
 \sigma_\pm=r^{2-\lambda_\gamma}V_\pm,\qquad
 \lambda_\gamma=\frac{12\gamma}{1+2\gamma},\qquad 0<c\le V_\pm\le C,
\end{equation}
and the complete stopped transition
\begin{equation}\label{sharp:upper-det-composition}
 x=x_b+\sigma_+\delta,\qquad
 \Pi(x)=y_b+\sigma_-O_\theta(D_\theta(I_\theta(\delta))),
 \quad I_\theta(\delta)=\xi_D,\quad 0<\delta\le\delta_*<1.
\end{equation}
The exact centres, cut factors, bridges and divided closing remain in
these maps. \(I,O\) vanish at zero, have derivative units, and are
ordinary \(C^3\) with bounded inverses. The incoming and outgoing saddle
chords have separate buffers. Their common power makes the quotient of
the two absolute gains a unit; the individual gains need no positive
lower bound.

The saddle of \eqref{deg:det-field} belongs to a compact analytic
family: its two rates have uniform positive moduli and
\(\wp\in[\wp_0,\wp_1]\Subset(0,\infty)\). Reduce the independent
incoming and outgoing chords separately to common compatible local
sections; their actual gap holonomies are ordinary \(C^3\) with
positive derivative units. Lemma~\ref{lem:uniform-positive-dulac}
therefore gives a local map \(cz^\wp e^{\mathcal R}\), with two
Euler remainder rows uniform through \(\wp=1\). Differentiation
replaces its remainder by
\(\Psi=\mathcal R+\log(1+z\mathcal R'/\wp)\).
The actual holonomies add \(O(z)+O(z^{\wp_0})\) to its value and
first Euler row. Thus, for fixed
\(0<\alpha<\min\{1,\wp_0\}\), the actual gap-to-gap map satisfies
\begin{equation}\label{sharp:upper-det-dulac}
 D_\theta'(z)=C_\theta\wp z^{\wp-1}e^{\Psi_\theta(z)},\qquad
 C_\theta^{\pm1}\le C,\qquad
 |\Psi_\theta|+|z\Psi_\theta'|\le Cz^\alpha(1+|\log z|).
\end{equation}
The estimate remains uniform at ratio one because it uses no division
by \(\wp-1\).

Put \(i_0=I'(0)>0\), \(o_0=O'(0)>0\). Differentiating the
complete composition in its original affine input gives exactly
\begin{equation}\label{sharp:upper-det-absolute}
 \begin{gathered}
 \Pi'(x_b+\sigma_+\delta)
 =\frac{\sigma_-}{\sigma_+}O'(D(I(\delta)))D'(I(\delta))I'(\delta)
 =\mathcal A_\theta\delta^{\wp-1}e^{R_\theta(\delta)},\\
 \mathcal A_\theta=\frac{\sigma_-}{\sigma_+}o_0C_\theta\wp i_0^\wp,
 \quad 0<c\le\mathcal A_\theta\le C,\qquad
 |R_\theta|+|\delta R_\theta'|\le C\delta^\alpha(1+|\log\delta|).
 \end{gathered}
\end{equation}
Here the full remainder is
\[
 R=\Psi(I(\delta))+(\wp-1)\log\frac{I(\delta)}{i_0\delta}
       +\log\frac{I'(\delta)}{i_0}
       +\log\frac{O'(D(I(\delta)))}{o_0}.
\]
The ordinary \(C^3\) bounds give \(I/(i_0\delta)=1+O(\delta)\)
in the required Euler rows, while \(D(I)\) and its first Euler row
are \(O(\delta^{\wp_0})\). Thus every term has the stated bound.
The common power in \eqref{sharp:upper-det-scales} cancels exactly,
without restricting \(\gamma\log r\); at \(\wp=1\) this full
remainder is still used.

Finally put \(S_D=r^2(1+p^2)\). The actual set cover is
\begin{equation}\label{deg:det-height-cover}
 \begin{gathered}
 \mathcal I_{\rm lo}=\{x\in\mathcal I_\theta:H_\theta\le2K_HS_D\},\qquad
 \mathcal I_{\rm hi}=\{x\in\mathcal I_\theta:H_\theta>K_HS_D\},\\
 \mathcal I_\theta=\mathcal I_{\rm lo}\cup\mathcal I_{\rm hi},\qquad
 \mathcal I_{\rm lo}\subset\{0<\xi_D<\xi_*\}\cup
 \{\xi_D\ge\xi_*/2\text{ in compact strict receivers}\}.
 \end{gathered}
\end{equation}
The low bound \(Y(0)\le2K_H(1+P_D^2)\) gives the finite
strict cover above; shrink \(\xi_*\) so its stopped part lies in
\eqref{sharp:upper-det-composition}. On strict cores the finite-cut
variations are units, and the common raw power
\eqref{sharp:upper-det-scales} cancels before the actual bridges and
divided closing bound \(\Pi'\). The signed high-height corollary
bounds the same derivative above.

The physical overlaps are \(\xi_*/2<\xi_D<\xi_*\) and
\(K_HS_D<H_\theta<2K_HS_D\). Gap equality faces belong to the
adjoining receiver, and both height equality faces are low. Window,
speed and cap clips preserve the cover and all intervening inputs;
a removed overlap places the connected retained interval in one
adjoining receiver. All these choices are made before selecting a
displacement root.
\end{proof}

\subsubsection{The stopped trace passage}
\label{subsec:trace-passage}

The trace chart contains a central saddle--node and an additional
equatorial saddle in the same shrinking neighbourhood. The central
normal form alone does not describe the return: its outgoing height
must also pass through this saddle. We include both passages in the
order of Figure~\ref{fig:trace-extra-saddle}.

\begin{proposition}[Uniform stopped-trace estimate]
\label{deg:prop-stopped-trace}
Fix \(0<p_0<p_1<\infty\), the passive compact coefficient block,
and the original buffered raw windows. There are independent bounds
\(r_*,\gamma_*>0\), an opening bound \(a_*>0\), a fixed positive
unscaled raw-height cap \(h_*\), and a final multiplier bound \(m_*\)
such that, for
\[
 \alpha_0=r^2a_L,\qquad \alpha_1=rp_L,\qquad
 0<r<r_*,\quad |\gamma|<\gamma_*,\quad |a_L|\le a_*,\quad
 p_0\le|p_L|\le p_1,
\]
the single original receiver
\begin{equation}\label{deg:trace-actual-receiver}
 \mathcal I_\theta=\mathcal I_\theta(h_*)
 =\{x:a_\theta+x\in J_\theta^{\rm out},\quad
       0<H_\theta(x)<h_*\},\qquad H_\theta(x)=z_\theta(\gamma;a_\theta+x),
\end{equation}
has \(F'>1/2\) everywhere for \(p_L<0\), and
\(F'=0\Rightarrow F''>0\) for \(p_L>0\), where
\(F(x)=x-m\Pi(x)\), \(0<m\le m_*\).
It consequently has at most one or two isolated zeros, respectively,
counted with total multiplicity. These conclusions apply to every
retained interval with the original window, speed and height clips,
including all intervening nonroot inputs and artificial equality
dividers. Both signs of \(a_L\), \(a_L=0\), the closed trace endpoints
and empty receivers are included; natural stopped endpoints are excluded.
The positive trace has a concave terminal receiver overlapping its
compact strict complement. No relation among \(m,r,\gamma\) is imposed.
\end{proposition}

% Figure contract: the stopped trace and its actual return composition.
% Native LaTeX picture; schematic, not a numerical phase portrait.
\begin{figure}[htb]
\centering
\setlength{\unitlength}{1pt}
\begin{picture}(430,238)
  % Panel (a): the outgoing extra saddle for positive central trace.
  \put(108,228){\makebox(0,0){\small (a) Positive-time geometry in $(q,Y)$}}
  \put(18,108){\line(1,0){198}}
  \put(216,108){\makebox(0,0)[l]{\small $q=-X$}}
  \put(38,108){\line(0,1){91}}
  \put(38,202){\makebox(0,0)[r]{\small $\Sigma_Q$}}
  \put(198,108){\line(0,1){91}}
  \put(198,202){\makebox(0,0)[l]{\small $\Sigma_-$}}
  \put(38,96){\makebox(0,0){\small $Q$}}
  \put(198,96){\makebox(0,0){\small $q_H$}}
  % Equator and its broken limiting connection.
  \linethickness{.35pt}
  \put(55,108){\line(1,0){95}}
  \put(150,108){\line(1,0){48}}
  \put(101,115){\makebox(0,0){\scriptsize equator $Y=0$}}
  \put(150,108){\circle*{5}}
  \put(150,95){\makebox(0,0){\scriptsize $e_R$}}
  \put(150,83){\makebox(0,0){\scriptsize extra saddle}}
  % Boundary connection to the saddle and its unstable branch.
  \linethickness{.8pt}
  \put(55,108){\vector(1,0){78}}
  \qbezier(150,108)(162,119)(174,132)
  \qbezier(174,132)(186,146)(198,163)
  \put(168,125){\vector(1,1){10}}
  \put(163,143){\makebox(0,0)[r]{\scriptsize unstable branch}}
  % Both curves increase with q. The strict orbit stays above the
  % unstable separatrix throughout their shared q interval.
  \linethickness{.55pt}
  \qbezier(38,154)(73,155)(108,158)
  \qbezier(108,158)(138,161)(162,171)
  \qbezier(162,171)(182,180)(198,193)
  \put(86,156){\vector(1,0){12}}
  \put(165,172){\vector(2,1){12}}
  \put(104,172){\makebox(0,0){\scriptsize typical strict orbit}}
  \put(38,183){\makebox(0,0)[l]{\scriptsize trace on $\Sigma_Q$}}
  \put(198,214){\makebox(0,0)[r]{\scriptsize fixed raw cut}}
  \put(108,66){\makebox(0,0){\scriptsize the axis segment ends at $e_R$;}}
  \put(108,54){\makebox(0,0){\scriptsize the positive-height orbit passes above it}}
  % Panel (b): actual composition used for the displacement.
  \put(325,228){\makebox(0,0){\small (b) The actual return composition}}
  \put(246,195){\makebox(0,0){\small weak input $t$}}
  \put(246,129){\makebox(0,0){\small weak input $t$}}
  \put(406,195){\makebox(0,0){\small affine input $x$}}
  \put(406,129){\makebox(0,0){\small raw output}}
  \put(278,195){\vector(1,0){30}}
  \put(340,195){\vector(1,0){30}}
  \put(278,129){\vector(1,0){30}}
  \put(340,129){\vector(1,0){30}}
  \put(307,207){\makebox(0,0){\scriptsize $M_d$ (DIR)}}
  \put(371,207){\makebox(0,0){\scriptsize $A_\theta$}}
  \put(307,141){\makebox(0,0){\scriptsize $\mathfrak H_\theta(t)=b_\theta+B_\theta(t)$}}
  \put(371,141){\makebox(0,0){\scriptsize $O_\theta$}}
  \put(311,180){\makebox(0,0){\scriptsize central passage}}
  \put(311,113){\makebox(0,0){\scriptsize actual outer Dulac}}
  \put(325,94){\makebox(0,0){\scriptsize includes $e_R$ and the fixed-section maps}}
  \put(325,78){\makebox(0,0){\scriptsize $\Pi\bigl(A_\theta(M_d(t))\bigr)=O_\theta(\mathfrak H_\theta(t))$}}
\end{picture}
\caption{The stopped trace and the map used in its zero estimate.  In
panel~(a), $q=-X$ increases in positive time.  The broken limiting
connection consists of the equatorial segment from the trace section
$\Sigma_Q$ to the additional equatorial saddle $e_R$, followed by its
unstable upper-half-plane branch toward the fixed raw section $\Sigma_-$.
The thinner curve is a representative strict positive-height orbit;
it remains above the unstable branch on their common $q$-interval.
The label \emph{trace on \(\Sigma_Q\)} marks the incoming height supplied by
the central saddle--node passage.  All distances and slopes are schematic.
Panel~(b) records the two sides of the physical composition: the central
saddle--node passage is the inverse Dumortier--Ilyashenko--Rousseau (DIR)
normal-form map $M_d$, while the outgoing height has the retained offset
$b_\theta$ and is then sent through the actual outer Dulac map $O_\theta$,
which includes the additional saddle and the fixed-section connectors.
Arrows in panel~(b) denote maps, not flow trajectories.}
\label{fig:trace-extra-saddle}
\end{figure}


\begin{proof}
For positive trace, we first estimate the outgoing equatorial passage
in its actual height coordinate. We then join it to the inverse central
passage and the incoming section map. Their logarithmic derivatives,
together with the regular connectors, give concavity near the stopped
orbit. A compact strict comparison and the high-height estimate extend
the count over the whole receiver. For negative trace, a direct
first-variation identity suffices.

Put \(v=\gamma+rX\), \(z=r^2Y\), and
\(\ell(v)=1-Av+Cv^2\), \(q_\gamma=1+2\gamma\). A positive time change gives
\begin{equation}\label{deg:trace-field}
\begin{split}
 X'&=a_L+p_LX-3\gamma X^2-rX^3-\ell(\gamma+rX)Y,\\
 Y'&=-Y\{(1+2\gamma)X+rX^2+Cr(\gamma+rX)Y\}.
\end{split}
\end{equation}
Keep the original affine \(x=h-a_\theta\) and
\(\Pi=\mathcal C_{\theta,m}\circ T_\theta(a_\theta+\cdot)\).
Lemma~\ref{phys:lem-c3-adapters} bounds divided closing through
ordinary order three, including \(m=0\). Phase derivatives fix all
coefficients and the actual multiplier. The two window admissions,
height and speed clips in \eqref{loc:nested-buffers} make
\eqref{deg:trace-actual-receiver} an interval before capture.

Choose the finite normal-form order and an initial \(|a_L|\le a_0\).
Take the signed high-height constant, increase \(K_H\) so that
\(|\bar a|\le a_0/K_H\), \(|\bar p|\le p_1/\sqrt{K_H}\),
\(\bar Y(0)=1\) lies inside its compact source block, and put
\(K_L=2K_H\). Fix \(X=L_c\) beyond the central equilibria.
The barrier estimate \eqref{sharp:upper-trace-barrier} below fixes a
bounded actual gap range before the terminal chords are selected.
Retain each auxiliary cut's exact continuation, centre and gain to
the original port. First take \(p_L\in[p_0,p_1]\).

\paragraph{The outgoing equatorial saddle.}
Set \(q=-X\).  The exact equations are
\begin{equation}\label{deg:trace-negative-field}
\begin{aligned}
 q'&=M(q)+L_-(q)Y,&M(q)&=-a_L+p_Lq+3\gamma q^2-rq^3,\\
 Y'&=Y\{qB(q)+\zeta(q)Y\},&B(q)&=1+2\gamma-rq,\\
 L_-(q)&=\ell(\gamma-rq),&\zeta(q)&=Cr(rq-\gamma).
\end{aligned}
\end{equation}
Fix \(w_0=1/8\), \(w_1=3/4\), and an auxiliary raw width
\(\delta\) inside the original ports, sufficiently small for this
fixed band. Choose a small \(a_0\) and a large \(Q\); then shrink
\(r,|\gamma|\), initially with \(|a_L|\le a_0\).
There is exactly one root \(R>Q\) of \(M\): \(M(Q)>0\),
\(M'(Q)>0\), and \(M'\) has exactly one positive root, beyond
which \(M\) decreases to \(-\infty\).  At \(R\),
\[
 rR^2=p_L+3\gamma R-a_L/R,\qquad
 \kappa=-M'(R)=p_L+rR^2-2a_L/R\ge p_0/2.
\]
Evaluating \(M\) on fixed compact intervals shows that \(R\to\infty\).
The identity \(rR=3\gamma+p_L/R-a_L/R^2\) then gives \(rR\to0\),
uniformly in the retained parameters.  The tangent and normal
eigenvalues are \(-\kappa\) and \(RB(R)>0\).  Thus the root is a
saddle, its incoming equatorial branch points from \(Q\) towards \(R\),
and its characteristic number satisfies
\[
 \beta_R=\frac{\kappa}{RB(R)}\longrightarrow0.
\]

Put
\[
 \begin{gathered}
 h(q)=\frac{M(q)}{q^2B(q)},\quad b(q)=qh(q),\quad
 \mathcal A(s)=\int_Q^s\frac{dq}{b(q)},\\
 \mathcal D=b(s)\partial_s,\qquad
 \epsilon=|a_L|/Q^2+p_1/Q+|\gamma|+rR.
 \end{gathered}
\]
This is the outgoing construction of
Lemmas~\ref{ann:lem-global-action-preparation} and
\ref{ann:lem-negative-gate}, with
\[
 (s,p,1)\quad\hbox{in their axis polynomial replaced by}\quad
 (r,p_L,-a_L).
\]
We verify the changed inputs. The first five Euler derivatives of
\(h=(-a_L/q^2+p_L/q+3\gamma-rq)/B(q)\) are \(O(\epsilon)\)
on \(Q/2\le q\le2R\) and its fixed logarithmic collars; the same
holds for \(L_-/B-1\) and \(\zeta q/B\). Moreover
\(b>0\) up to \(R\), \(|b'|\le C\epsilon\), and
\(\mathcal A\to\infty\) at \(R\). On the fixed entrance interval,
\[
 M(q)/q\ge p_0-2a_0/Q-3|\gamma|Q-rQ^2\ge p_0/2
 \qquad(Q/2\le q\le Q).
\]
After choosing \(R_0>2Q\), the bound \(b\le Cp_1\) on
\([Q,R_0]\) follows by an independent shrink of \(r,|\gamma|\).

For clarity, the two equations used in that construction are
\begin{equation}\label{deg:trace-normal-field}
 \begin{gathered}
 u'=h+cw,\qquad w'=w\{1-2h+(k-2c)w\},\\
 u=\log q,\quad w=Y/q^2,\quad c=L_-/B,\quad k=\zeta q/B,
 \end{gathered}
\end{equation}
in positive normal time \(d\tau_n=qB(q)\,dT\), and
\begin{equation}\label{deg:trace-graph}
 \frac{dw}{d\log q}=F(\log q,w)
 =w\left\{\frac{1+2\gamma-rq+Crq(rq-\gamma)w}
 {L_-(q)w+3\gamma-rq+p_L/q-a_L/q^2}-2\right\}.
\end{equation}
The first is unchanged in form, so
Lemmas~\ref{lem:normal-time-projection} and
\ref{lem:normal-time-action-height} give the strong footpoint and
its actual fixed-cut height on \(R_0\le s<R\):
\begin{equation}\label{deg:trace-height}
 \log y=-\mathcal A(s)+2\log s+\log U(s)+E_Q(y),\qquad
 0<y\le e_0=C_QR_0^2e^{-c(R_0-Q)/p_1}.
\end{equation}
Here \(U\) and its reciprocal and logarithmic derivatives are bounded,
\(E_Q(0)=0\), \(\|E_Q\|_{C^3}\le C_Q\), and
\begin{equation}\label{deg:trace-height-derivative}
 \mathcal D\log y=-K,\qquad
 |K-1|+|\mathcal DK|\le C\epsilon+C_Qe_0,\qquad K\ge\tfrac12.
\end{equation}
The passage stays in \(Q/2<q<2R\), and \(y\downarrow0\) at \(R\).

For the continuation, fix \(|\gamma|<\delta/2\). The strong-section
abscissa is \(q_s(s)=sU_s(s)\), with \(q_s'\asymp1\) and
three bounded logarithmic derivatives of \(U_s\). On
\(w_0\le w\le w_1\), the denominator in
\eqref{deg:trace-graph} is bounded below,
\(F(w_0)>c\), \(F(w_1)<-c\), \(F_w\le-1\), and its required
state and logarithmic-radius derivatives are bounded. Since
\(q_s(R)<q_H=(\delta+\gamma)/r\), the same invariant-strip and
moving-entrance argument as in
\eqref{ann:repair-moving-variation}--\eqref{ann:repair-completed-log-jets}
gives a unique raw exit and bounded logarithmic phase derivatives.
Only their first row is needed here. The new term
\(-a_L/q^2\) has the same integrable weight as the former \(q^{-2}\),
with bounded coefficient; no reciprocal \(\beta_R\) is introduced.

The full raw image is admitted before composing the closing. Indeed,
choose \(0<d_{\rm r}\ll\delta\), then require \(2rR<d_{\rm r}\).
Comparison from \(rq=d_{\rm r}\) to the raw exit with
\(w_\Gamma(\xi)=(\sqrt{1+\xi^2}+1)^{-1}\) gives, uniformly in \(s\),
\[
 |W_H(s)-w_\Gamma(\delta+\gamma)|
 \le C d_{\rm r}/\delta+
 C_{d_{\rm r},\delta}(|A-1|+|C-1|+|\gamma|+rp_L+|a_L|r^2).
\]
This is the raw comparison in the negative-gate proof with the same
changed polynomial. Choose \(d_{\rm r}/\delta\) from the original
window reserve, then shrink the coefficients; the raw height is
\((\delta+\gamma)^2W_H\), with organizing-height translation
\(O(\delta|\gamma|)\). The compact bridge and divided closing are
therefore defined on the whole image. Their ordinary \(C^3\) bounds
give the completed output \(\Gamma_-(s)\) with
\[
 \Gamma_-'<0,\qquad
 |s\partial_s\log(-\Gamma_-')|\le C,\qquad
 |\mathcal D\log(-\Gamma_-')|\le C\epsilon.
\]
The two physical radii \((\delta\pm\gamma)/r\) remain distinct in
their absolute gains.

Define \(O(y(s))=\Gamma_-(s)\). The same logarithmic chain rule as
in Proposition~\ref{ann:prop-endpoint-curvature} now gives
\begin{equation}\label{deg:trace-extra-saddle}
 \begin{split}
 O'(y)&=\frac{(-\Gamma_-')b}{yK}>0,\\
 y\frac{O''}{O'}
 &=-1-\frac{\mathcal D\log(-\Gamma_-')+b'-\mathcal D\log K}{K}
 =-1+O(\epsilon+C_Qe_0)<-\tfrac12.
 \end{split}
\end{equation}
It holds on the actual height interval \(0<y<y(R_0)\), uniformly
as \(\beta_R\to0\); no derivative at the excluded endpoint is used.

\paragraph{The central passage and its connectors.}
At \(a_L=r=\gamma=0\) the central eigenvalues are \(p_L,0\),
with the nonzero weak quadratic coefficient of
Lemma~\ref{deg:lem-positive-atlas}. Apply the finite-smooth
multiparameter orbital normal form of
\cite[Theorem~3 and formula~(2.2), October 2000 author copy p.~6]{DIR02}
to the reversed field, at the preselected joint order \(N\ge12\).
Its positive orbital factor gives the inverse physical passage, with
no scalar-opening rank condition.

The axis branch has \(X_{\rm ax}=-a_L/p_L+O(a_L^2)\), \(Y=0\),
and the interior branch satisfies
\[
 \left.\partial_{a_L}Y_{\rm int}\right|_{a_L=0}
 =\frac{1+2\gamma}
 {\ell(\gamma)(1+2\gamma)+p_LCr\gamma}>c>0.
\]
Their weak separation is \(a_L\) times a positive unit; the strong
rate stays nonzero. A finite compact cover gives common boxes and phase
bounds across \(a_L=0\). The physical boundary is the interior saddle
for \(a_L>0\), the axis saddle for \(a_L<0\), and their common root
at zero. Center this weak root at \(x_+(d)\), with the other at
\(x_+(d)-d\), \(d\asymp|a_L|\). Only signed-side phase derivatives
are needed.

In reversed normal-form time take \(x'=g_d(x)\), \(\eta'=-\eta\),
a fixed strong entrance \(\eta=\eta_0>0\), and a fixed weak exit
\(x=\rho\). The orbit from \((x_+(d)+t,\eta_0)\) reaches it at
height \(\eta_0M_d(t)\). With bounded normal-form coefficient
\(\mathfrak a\), the two weak factors give
\begin{equation}\label{deg:trace-stopped-map}
 \begin{gathered}
 G_d(t)=\frac{t(t+d)}{1+\mathfrak a(x_+(d)+t)},\qquad
 g_d(x_+(d)+t)=G_d(t),\\
 M_d(t)=\exp\!\left[-\int_{x_+(d)+t}^{\rho}\frac{dx}{g_d(x)}\right].
 \end{gathered}
\end{equation}
For small \(d\ge0\), \(t>0\),
\begin{equation}\label{deg:trace-stopped}
 M_d'=\frac{M_d}{G_d},\qquad
 \frac{M_d''}{M_d'}=\frac{1-G_d'}{G_d},\qquad
 1-G_d'\ge\delta_G>0,\quad G_d\asymp t(t+d).
\end{equation}
The quantities \(G_d\) and \(M_d\) are uniformly small on a short
terminal interval: \(G_d\le Ct(t+d_*)\) for \(0\le d\le d_*\),
while the lower action bound \(c\int_t^{\rho/2}du/[u(u+d_*)]\),
with fixed small \(d_*>0\), diverges at zero.
Keeping the two coalescing factors together makes this estimate
uniform as their separation vanishes.

Absorb the fixed factor \(\eta_0\) into the incoming section map.
From the central strong exit to \(q=Q\), the flight has fixed
finite length and \(q'\ge cq>0\).  Ordinary graph variations give
\begin{equation}\label{deg:trace-height-connector}
 \begin{gathered}
 \mathfrak H_\theta(t)=b_\theta+B_\theta(t),\qquad B_\theta(0)=0,\\
 0<c_Q\le B_\theta'\le C_Q,\qquad
 |\mathfrak H_\theta''/\mathfrak H_\theta'|\le C_Q,\qquad
 0\le b_\theta\le C_Q|a_L|.
 \end{gathered}
\end{equation}
Here \(b_\theta=0\) for \(a_L\le0\). For \(a_L>0\), the
interior saddle's strong branch has a uniform spectral gap; its
smooth finite flight to \(Q\) gives the positive offset. Retain this
offset in \(O\), using phase regularity on each signed side.

After \(Q,R_0\) are fixed, \(y_\theta(R_0)\) has a uniform
positive lower bound.  At \(r=\gamma=0\) it is the height of a
finite regular passage on a compact parameter block, and continuity
preserves positivity.  Choose \(y_*>0\) below this bound.  Only
now shrink \(a_*,t_*\) so that
\begin{equation}\label{deg:trace-domain-inclusion}
 \mathfrak H_\theta([0,2t_*])<y_*/2<y_\theta(R_0).
\end{equation}
Thus every strict terminal orbit, including those with positive
offset, lies in the domain of \eqref{deg:trace-extra-saddle}.

For the incoming connector, at \(r=\gamma=0\) the boundary branch
beyond the positive local cut satisfies \(X,Y>0\) and
\(D=Y-a_L-p_LX>0\).  At \(D=0\), reverse time points into
\(D>0\), since \(\dot D=p_LD-XY\).  Also \(Y_X\le2X\)
for large \(Y\) on bounded \(X\)-intervals.  These facts continue
the branch to every fixed positive cut without loss of its graph
or finite-\(X\) blow-up.  At large fixed \(X\),
\[
 \frac{dW}{d\log X}=\frac{W}{W-p_L/X-a_L/X^2}-2W,
 \qquad W=Y/X^2.
\]
Choose the cut so that \(0<p_L/X+a_L/X^2<1/4\).  The right
side is at least \(1/2\) for \(0<W\le1/4\) with \(D>0\),
and is negative for \(W\ge1\).  The compact family enters a
common quadratic corridor at a fixed cut \(L_t\), as does a small
chord of strict incoming orbits for nearby coefficients.
For nearby coefficients the rational tail equation has
\(F_W\le-1\) and bounded \(F_{WW}\). Its first variation decays
exponentially, so the relative-jet argument
\eqref{tail:relative-jet-equations} bounds the pre-Schwarzian uniformly.
The same fixed raw comparison as on the negative side admits the whole
chord into the positive window. After the finite connector and original
bridges, the actual input adapter therefore satisfies
\begin{equation}\label{deg:trace-input-adapter}
 x=A_\theta(M_d(t)),\qquad A_\theta'>0,\qquad
 |A_\theta''/A_\theta'|\le C_Q.
\end{equation}
Its absolute first derivative needs no lower bound.

\paragraph{Curvature and the overlap of actual phase intervals.}
The three constructed maps refer to the same orbit by flow uniqueness:
\(A_\theta\) continues to the actual input, \(\mathfrak H_\theta\)
to \(q=Q\), and \(O_\theta\) through the extra saddle and closing.
Thus
\[
 x=A_\theta(M_d(t)),\qquad \Pi(x)=O_\theta(\mathfrak H_\theta(t)).
\]
All factors are increasing. Put \(H=\mathfrak H_\theta(t)\).
Differentiation in \(t\), keeping the
affine input \(x\) in the displacement, gives
\begin{equation}\label{deg:trace-composition}
 \begin{split}
 G_d\frac{d}{dt}\log\Pi'(x(t))
 ={}&G_d\frac{O''(H)}{O'(H)}H'+G_d\frac{H''}{H'}\\
 &-M_d\frac{A''(M_d)}{A'(M_d)}-(1-G_d').
 \end{split}
\end{equation}
The first term is nonpositive by \eqref{deg:trace-extra-saddle}
and \eqref{deg:trace-domain-inclusion}.  The next two are bounded
by \(C_Q(G_d+M_d)\), whereas the last is at most \(-\delta_G\).
Shorten the terminal interval so that the whole expression is at
most \(-\delta_G/2\).  Since \(G_d>0\) and \(x'(t)>0\),
this proves \(\Pi''<0\) in the original affine input.

Fix a positive local cut \(X=L_c\) beyond the nearby central
equilibria, and let \(Y_b^+(L_c)\) be the incoming boundary
height.  The normalized actual gap is
\begin{equation}\label{deg:trace-actual-gap}
 \xi=Y(L_c)-Y_b^+(L_c)
 =\frac{z(\gamma+rL_c)-r^2Y_b^+(L_c)}{r^2}.
\end{equation}
Thus \(r^2\xi\) is the actual raw gap on this section, with the
same zero, ordering and orientation as \(\xi\).
Order and the compact local connector make \(\xi(t)\) strictly
increasing. More precisely, the connector on the fixed local cuts
gives \(\xi(t)=\psi_\theta(M_d(t))-\psi_\theta(0)\), with
\(0<c\le\psi_\theta'\le C\). Since \(G_d(u)\ge cu^2\),
the action in \eqref{deg:trace-stopped-map} at \(t=t_*/2\) is
at most \(C/t_*\). Consequently
\(\xi_\theta(t_*/2)\ge c\exp(-C/t_*)>0\), uniformly on the
completed coefficient block, including \(d=0\). Choose
\begin{equation}\label{deg:trace-overlap}
 0<2\xi_0<\inf_\theta\xi_\theta(t_*/2).
\end{equation}
Then \(0<\xi\le2\xi_0\) lies in the terminal interval.
The whole lower range \(Y_0=H_\theta(x)/r^2\le K_L\) reaches this
same cut. Indeed its exact scalar graph is
\begin{equation}\label{sharp:upper-trace-barrier}
 Y_X=\frac{Y\{(1+2\gamma)X+rX^2+Cr(\gamma+rX)Y\}}
 {\ell(\gamma+rX)Y-a_L-p_LX+3\gamma X^2+rX^3}.
\end{equation}
The through denominator is positive. Above a fixed \(M_0>K_L+1\)
it is at least \(\ell_0Y/2\), and on \(0\le X\le L_c\)
\(Y_X\le C_{L_c}+Y\). Comparison after the first crossing of \(M_0\)
and Gronwall give the previously fixed bound
\begin{equation}\label{sharp:upper-trace-bounded-gap}
 0<\xi=Y(L_c)-Y_b^+(L_c)\le B.
\end{equation}
Only original through admission is used here. The divider height is
nonnegative and is subtracted on the same section.

On \(\xi\ge\xi_0\) the limiting passage is uniformly separated from
its stopped graph. To obtain this separation before constructing a strict
cover, let \(t_{\rm ref}\) be the unique phase with \(\xi=\xi_0\).
The uniform vanishing of \(M_d(t)\) bounds \(t_{\rm ref}\) away from
zero. The standing normal form gives a nonbroken physical passage at
every such positive phase: for \(a_L>0\) it crosses above the interior
saddle with \(Y_0>a_L\), and for \(a_L\le0\) with \(Y_0>0\).
Take the positive minimum over the compact reference-phase and
coefficient block at \(r=\gamma=0\); scalar order then gives
\[
 Y_0-\max\{a_L,0\}\ge\eta_0>0\qquad(\xi\ge\xi_0).
\]
For \(D=Y-a_L-p_LX\), \(D(0)>0\) and
\(\dot D=p_LD-XY\) excludes a first zero in either relevant half-axis
orientation. The bound \(Y_X\le C+Y\) excludes finite-\(X\) blow-up.
Thus the bounded complement has a compact strict-family cover with
positive core height and speed margins. At large \(|X|\), its
\(W=Y/X^2\) equation points upward below \(1/4\) and downward above
\(1\), so both half-tails enter quadratic corridors at a common fixed
cut. Their actual outward continuations and fixed raw comparison give
both original-window margins. In particular
\(\xi_0\le\xi\le2\xi_0\) is a physical terminal/strict overlap.

For a strict limiting graph, Proposition~\ref{ann:prop-source-curvature}
with \(a=a_L\), \(p=p_L\), \(b=B_+\), and
\(\sigma=\partial_bY>0\) gives
\begin{equation}\label{deg:trace-strict-curvature}
 \frac{S''}{S'}=-2\int_{\mathbb R}\frac{X(a_L+p_LX)\sigma}{D^3}\,dX
              =-p_L\int_{\mathbb R}\frac{\sigma}{D^2}\,dX<0.
\end{equation}
The compact strict cover bounds \(\sigma\) below and \(D\) above
on a common core; \(p_L\ge p_0\) gives a uniform negative margin.

To transfer this margin, use the exact normalization
\eqref{deg:det-normalized-tail}. The changed coefficient \(a_L\)
gives the weight
\[
 \|N\|_{C^3}\le C\{rR^{-2+\nu_0}+|\gamma|R^{-3+\nu_0}
                  +|p_L|R^{-4+\nu_0}+|a_L|R^{-5+\nu_0}\},
 \qquad 0<\nu_0<1.
\]
The argument of Proposition~\ref{prop:source-normalization}, with this
integrable weight and the admitted enlarged chords, gives phase \(C^3\)
convergence independently in \(r,\gamma\). Exterior factors normalized
by their exact value and slope converge to the identity by
\eqref{deg:det-adapter-normalization}--\eqref{deg:det-divided-closing}.
Use the actual maps \(h(b),k(b)\) and their normalization
\eqref{eq:reference-normalization}, with
\(c_h=h'(b_*)>0\), \(c_k=k'(b_*)>0\). Let
\(\beta=b_*+(x+a_\theta-h(b_*))/c_h\). As in
\eqref{ann:p0-two-actual-adapters},
\begin{equation}\label{deg:trace-physical-curvature}
 \Pi''(x)=\frac{c_k}{c_h^2}\Phi''(\beta)<0.
\end{equation}
Both raw first variations are \(r^{2-\lambda_\gamma}\) times
positive units. Their quotient, the actual bridges and divided closing
bound \(\Pi'\) on the strict pieces and their regular finite-flight
ends; the individual gains may vanish.

At every actual point with \(Y_0>K_H\), freeze the dilation at that
point before a phase derivative. The exact normalized parameters are
\begin{equation}\label{sharp:upper-trace-high-dilation}
 (\bar a,\bar p,\bar\gamma,\bar s)=
 (a_L/Y_0,p_L/\sqrt{Y_0},\gamma,r\sqrt{Y_0}),\qquad \bar Y(0)=1.
\end{equation}
The fixed unscaled cap \(h_*\) gives \(\bar s<\sqrt{h_*}\).
The signed compact source block of
Corollary~\ref{rew:cor-high-derivative-signed} receives these parameters.
Its two absolute original-port variations have the same
\(\bar s^{2-\lambda_\gamma}\) power, including every bridge and
\eqref{deg:det-divided-closing}. Their quotient bounds the derivative
of this same physical \(\Pi\) at every high input, including nonroots.

Thus the original set cover is
\[
 \mathcal I_\theta=\{Y_0\le K_L\}\cup\{Y_0>K_H\},\qquad
 \{Y_0\le K_L\}\subset\{0<\xi\le2\xi_0\}\cup\{\xi\ge\xi_0\}.
\]
Each set is restricted to \(\mathcal I_\theta\). The physical overlaps
are \(K_H<Y_0<K_L\) and \(\xi_0<\xi<2\xi_0\). Enlarged chords
receive both positive gap equality faces; the height equality faces are
low. Original window, speed or height clipping preserves this cover,
including all intervening inputs, and excludes the broken \(t=0\).
On the terminal and strict pieces \(\Pi''<0\); on the high and
regular pieces \(\Pi'\le C\). Choose \(m_*\) last so that the latter
have \(F'>1/2\). Every zero of \(F'\) has \(F''=-m\Pi''>0\).
Lemma~\ref{sharp:upper-critical-count} gives the whole positive-trace
bound.

\paragraph{The negative trace.}
Now take \(\alpha_1=-rp\), \(p\in[p_0,p_1]\), and
\(\alpha_0=r^2a_L\), \(|a_L|\le a_*\). Let \(\delta_{\rm raw}\)
bound \(|v-\gamma|\) on the original small raw chart. In the raw coordinate
\(u=v-\gamma\), put
\[
 \begin{gathered}
 \mathcal M=-r^2a_L+rpu+3\gamma u^2+u^3,\qquad
 \mathscr D=\mathcal M+\ell(v)z>0,\\
 \nu=\frac{6\gamma}{1+2\gamma},\qquad N=rp+(3-\nu)u^2\ge rp+2u^2.
 \end{gathered}
\]
The graph equation
\(z_v=f=z\{(1+2\gamma)u+u^2+Cvz\}/\mathscr D\)
satisfies \eqref{deg:raw-log-jacobian} with these \(\mathcal M,N\).
\begin{equation}\label{deg:trace-jacobian}
 f_z=\frac d{dv}\log\frac{z^{1+\nu}}{\mathscr D}
   +\frac{N+\{\ell'(v)+(1-\nu)Cv\}z}{\mathscr D}.
\end{equation}
On \(\mathcal M<0\),
\[
 |\mathcal M|\le a_*r^2+rp|u|+3|\gamma|u^2+|u|^3
 \le C(a_*r+\delta_{\rm raw}+|\gamma|)N.
\]
Let \(\kappa_0=\tfrac12\inf\ell>0\).  Where
\(\mathscr D\ge\kappa_0z\), the last quotient in
\eqref{deg:trace-jacobian} is uniformly bounded below.
On the complement, \(\mathcal M<0\) and
\(z\le|\mathcal M|/(\ell-\kappa_0)\); the displayed estimate
makes its numerator at least \(N/2\) on the standing small raw chart
after the independent coefficient shrink.  It is therefore nonnegative there.
Integrating from the positive raw section to the negative one
reverses the orientation of the interval and gives
\begin{equation}\label{deg:trace-negative-derivative}
 0<T'\le C\frac{(z^{1+\nu}/\mathscr D)_-}
 {(z^{1+\nu}/\mathscr D)_+}\le C_0.
\end{equation}
The endpoint ratio is bounded on the fixed actually admitted raw
windows: endpoint heights are positive units and graph speeds have
fixed positive lower bounds. No lower central speed is used. The
estimate holds on the full receiver \eqref{deg:trace-actual-receiver},
including its speed clip and every intervening nonroot input.
The ordinary divided closing has bounded positive derivative, so
\(0<\Pi'\le C_{\rm neg}\). Add \(m_*C_{\rm neg}<1/2\) to the final
multiplier choice; then \(F'>1/2\) throughout, giving at most one
simple zero. The positive-phase estimates remain uniform at \(a_L=d=0\).
\end{proof}


\begin{lemma}[Nonnegative-opening domains]
\label{deg:lem-nonnegative-domain}
On every coefficient fibre with \(\alpha_0\ge0\), the coefficient
alternatives and phase intervals of
Lemma~\ref{deg:lem-positive-atlas} cover the connected captured
physical interval, with their regular continuations and all points
between captured roots. The compact strict and regular portions have
bounded complete raw derivative. The common high-height interval is
contractive after the multiplier shrink of
Corollary~\ref{rew:cor-high-derivative-signed}, including the
positive-height equality face \(\alpha_0=\alpha_1=0\).
The only remaining stopped portions are the determinant passage of
Lemma~\ref{deg:lem-determinant-preparation} and the signed trace
passages of Proposition~\ref{deg:prop-stopped-trace}.
\end{lemma}

\begin{proof}
For each whole coefficient fibre use the alternative of
Lemma~\ref{deg:lem-positive-atlas}. The actual determinant cover
\eqref{deg:det-height-cover} and signed trace receiver
\eqref{deg:trace-actual-receiver} retain its stopped, strict, regular
and high pieces, including the original clips and every intervening
input. Their two absolute tail gains have a common power, which cancels
before the complete raw derivative is bounded on the strict and regular
pieces. The signed high-height estimate bounds that same derivative
above, including when both opening coefficients vanish. Choosing the
finite multiplier last makes these bounded pieces contractive.
The gap-zero completion is excluded. The whole-fibre sharp counts are
Propositions~\ref{sharp:upper-determinant} and \ref{sharp:upper-trace};
Corollary~\ref{sharp:upper-local-coefficient-ranges} takes the maximum
across coefficient alternatives.
\end{proof}
